%% Arithmax Research: Congressional Trading Signals Paper
%% Publisher: Frankline Oyolo (Frankline@arithmax.com)

\documentclass[twocolumn,ieee]{Base_Latex_Template/arithmaxresearch}

\usepackage{amsmath}
\usepackage{algorithm}
\usepackage{amssymb}
\usepackage{url}
\usepackage{float}
\floatname{algorithm}{Algorithm}

\begin{document}

% Title
\title{Information Decay and Optimal Rebalancing Frequency for Congressional Trading Signals}

% Author information (explicit IEEE block for robust rendering)
\author{%
\IEEEauthorblockN{Frankline Misango Oyolo}\\
\IEEEauthorblockA{\textit{Arithmax Research} \\
Arithmax Quantitative Research \\
Email: frankline@arithmax.com}%
}

% Title page with Arithmax branding
\maketitle
\begin{center}
\vspace{-1em}
\arithmaxtitlelogo[4cm]
\vspace{0.5em}
\end{center}

\begin{abstract}
This paper investigates the commercialization of public Congressional trading disclosures as a systematic alpha source. We develop a framework for extracting, processing, and trading legislative information signals with explicit transaction-cost accounting. Our key empirical result is a robust cost-signal trade-off: weekly rebalancing yields 1,422.67\% annual turnover with 1.42\% annual cost drag, while monthly rebalancing yields 859.33\% annual turnover with 0.86\% annual cost drag. In our tested design space (weekly, monthly, quarterly, semi-annual), monthly rebalancing is the best risk-adjusted configuration, improving Sharpe ratio by 77.7\% (0.226 to 0.402) and annualized return by 86.6\% (4.16\% to 7.76\%) versus weekly. We separate prior art from novel contribution: prior literature establishes Congressional signal relevance and baseline portfolio/execution theory, while our contribution is operational optimization under disclosure-lag constraints. We document methodology, empirical results, mathematical framing, and practical deployment implications for satellite-strategy use.
\end{abstract}

% =====================================================================
\section{Introduction}
% =====================================================================

The efficient market hypothesis (EMH) posits that stock prices fully reflect all available public information \cite{fama1970efficient}. However, empirical evidence suggests significant market anomalies persist, particularly around identified information events. The puzzle intensifies when considering information that is both publicly available \textit{and} disclosed through government channels: Congressional stock trading disclosures.

Members of Congress and their families are required by the STOCK Act (2012) to publicly disclose their securities transactions within 45 days of execution. These filings represent a unique information source: systematic insider knowledge (representatives likely have superior information) combined with mandatory public disclosure (eliminating legal information asymmetry). Prior research has documented abnormal performance in Congressional portfolios \cite{ziobrowski2004abnormal, ziobrowski2011house}, suggesting potential alpha generation opportunities.

Yet exploiting this opportunity presents a fundamental challenge: \textbf{information decay} coupled with \textbf{transaction costs}. While Congressional trades are disclosed publicly, the 30-45 day lag means traders act on already-aged information. Paradoxically, frequent rebalancing to capture fresh signals generates transaction costs that may exceed the signal's remaining economic value.

This paper makes three primary contributions:

\begin{enumerate}
\item \textbf{Quantification of cost-signal trade-off}: We empirically show that weekly scheduling materially over-trades relative to monthly in our implementation, with annual turnover falling from 1,422.67\% to 859.33\% and cost drag from 1.42\% to 0.86\%.

\item \textbf{Information decay modeling}: We develop a framework for understanding how Congressional trading signals decay over time and establish the relationship between disclosure lag, rebalancing frequency, and net alpha.

\item \textbf{Operational optimization}: We provide evidence-based guidance for practitioners on optimal rebalancing frequency, position sizing, and capacity constraints for this strategy class.
\end{enumerate}

% =====================================================================
\section{Data Architecture and Signal Formulation}
% =====================================================================

\subsection{Data Sources and Integration}

Our data foundation comprises three integrated layers:

\begin{enumerate}
\item \textbf{Congressional Trading Disclosures}: Scraped from official HOUSE.GOV and SENATE.GOV filing databases via Quiver Quantitative, containing 107,386 unique transactions spanning 2014-2026. Each record includes transaction date, filing date, transaction type (buy/sell), security identifier, transaction value, and representative metadata.

\item \textbf{Market Microstructure}: Daily OHLCV data from Yahoo Finance (95.5\% coverage) with Alpaca fallback (4.6\% coverage) for 3,991 securities. Data spans 2014-03-27 with 3,077 trading days.

\item \textbf{Reference Data}: Market capitalization and sector classification from Yahoo Finance, volume data for market impact estimation, and historical volatility measurements.
\end{enumerate}

\subsection{Signal Processing Pipeline}

Let $\mathcal{D} = \{d_1, d_2, \ldots, d_n\}$ denote the set of Congressional disclosures. For each security $s$ and time $t$, we define the aggregated Congressional signal as:

$$\mathcal{S}_{s,t} = \sum_{d_i \in \mathcal{D}_t} w(d_i) \cdot v(d_i) \cdot \text{sign}(d_i)$$

where:
\begin{itemize}
\item $\mathcal{D}_t$ = Congressional trades within lookback window $[t - L, t]$ (default $L = 45$ days)
\item $w(d_i)$ = temporal weight, decreasing with age of disclosure
\item $v(d_i)$ = transaction value (normalized by average daily volume)
\item $\text{sign}(d_i)$ = buy (+1) or sell (-1)
\end{itemize}

We employ inverse-volatility weighting:

$$w(d_i) = \frac{1}{\sigma_s(d_i)} / \sum_{j} \frac{1}{\sigma_s(d_j)}$$

This produces signals with magnitude in $[0, 1]$, reflecting consensus Congressional positioning.

\subsection{Data Quality and Coverage}

Statistical coverage analysis:
\begin{itemize}
\item Congressional trades: 107,386 transactions, mean \$87K per transaction
\item Market data coverage: 85\% of S\&P 1500 universe
\item Common universe: 2,847 securities with both Congressional trades and complete OHLCV history
\item Backtest period: 2014-01-02 to 2026-03-27 (12.2 years, 3,077 trading days)
\end{itemize}

% =====================================================================
\section{Strategy Design and Methodology}
% =====================================================================

\subsection{Portfolio Construction}

Given signals $\mathcal{S}_{s,t}$ for universe $\mathcal{U}_t$, we construct portfolios using risk parity (inverse volatility) weighting:

$$w_{s,t} = \frac{\mathcal{S}_{s,t} / \sigma_s(t)}{\sum_{s' \in \mathcal{U}_t} \mathcal{S}_{s',t} / \sigma_{s'}(t)}$$

subject to constraints:
\begin{align}
\sum_s w_{s,t} &= 1 \quad \text{(fully invested)} \\
w_{s,t} &\leq w_{\max} \quad \text{(position limit, } w_{\max} = 0.10 \text{)} \\
|n_t| &\leq n_{\max} \quad \text{(holding limit, } n_{\max} = 50 \text{)}
\end{align}

\subsection{Transaction Cost Modeling}

We model transaction costs as:

$$TC_t = P_t \cdot TO_t \cdot (c_{\text{comm}} + c_{\text{slip}})$$

where:
\begin{itemize}
\item $P_t$ = portfolio value at rebalance $t$
\item $TO_t$ = turnover ratio (sum of absolute weight changes)
\item $c_{\text{comm}} = 5$ bps (commission)
\item $c_{\text{slip}} = 5-10$ bps (slippage, market-cap dependent)
\end{itemize}

Total cost rate: $c_{\text{total}} = 0.10\%$ per unit turnover.

Annual cost impact: $\text{Cost}_{\text{annual}} = TO_{\text{annual}} \cdot c_{\text{total}}$, where $TO_{\text{annual}}$ is the annualized turnover.

\subsection{Baseline Strategy (Weekly Rebalancing)}

Original implementation rebalances every 7 calendar days (52 rebalances/year):

\begin{itemize}
\item Number of rebalances: 638 (actual trading dates)
\item Average turnover per rebalance: 27.4\%
\item Annual turnover: 1,422.67\%
\item Annual cost drag: 1.42\%
\item Annualized return (gross): 4.16\%
\item Sharpe ratio: 0.226
\end{itemize}

Performance diagnosis: While signals show positive correlation with forward returns (54.6\% win rate) and reasonable magnitude (0.86 average signal strength), transaction costs dominate, creating severe alpha leakage.

% =====================================================================
\section{Optimization: Rebalancing Frequency Trade-off}
% =====================================================================

\subsection{Hypothesis and Experimental Design}

\textbf{Core hypothesis}: Information decay from Congressional disclosures is slow (30-45 day cycle) relative to execution friction. Therefore, reducing rebalancing frequency from weekly to monthly should:
\begin{enumerate}
\item Reduce transaction costs materially (1,422.67\% → 859.33\% annual turnover)
\item Maintain signal potency (Congressional lag unchanged)
\item Materially improve net risk-adjusted returns
\end{enumerate}

We test four rebalancing frequencies:
\begin{align}
f \in \{W, ME, Q, SA\} \quad \text{(weekly, monthly, quarterly, semi-annual)}
\end{align}

All other parameters held constant: $L = 45$ days lookback, $n_{\max} = 50$ positions, $w_{\max} = 10\%$.

\subsection{Mathematical Formulation of Information Decay}

Let $\alpha(t, \tau)$ denote the alpha available from a Congressional disclosure $\tau$ periods after execution. We model this as:

$$\alpha(\tau) = \alpha_0 \cdot e^{-\lambda \tau} \quad \text{(exponential decay)}$$

where:
\begin{itemize}
\item $\alpha_0$ = initial alpha at disclosure ($\tau = 0$)
\item $\lambda$ = decay rate parameter
\item $\tau$ = time since Congressional trade execution
\end{itemize}

For Congressional trades with 45-day disclosure lag, we estimate $\alpha(\tau = 30-45)$ remains at 70-80\% of initial value. Conversely, transaction costs accumulate \textit{per rebalance}, not per unit time.

Define expected cost-adjusted return for rebalancing frequency $f$:

$$R_f^{\text{net}} = \sum_{t=1}^{T} r_t^{gross} - \sum_{i=1}^{n_f} TC_i$$

where $n_f$ = number of rebalances at frequency $f$ over period $[1, T]$.

Monthly rebalancing ($f = ME$) yields (in this sample):
\begin{align}
n_{ME} &\approx 11.9\ \text{rebalances/year}, \\
n_{W} &\approx 52.2\ \text{rebalances/year}, \\
TO_{ME}^{\text{annual}} &= 859.33\%, \\
TO_{W}^{\text{annual}} &= 1422.67\%, \\
\text{Cost}_{ME}^{\text{annual}} &= 0.0086, \\
\text{Cost}_{W}^{\text{annual}} &= 0.0142.
\end{align}

\subsection{Optimization Results Summary}

\begin{table}[h]
\centering
\scriptsize
\setlength{\tabcolsep}{3.5pt}
\renewcommand{\arraystretch}{0.95}
\begin{tabular}{lrrrr}
\hline
\textbf{Metric} & \textbf{Weekly} & \textbf{Monthly} & \textbf{Quarterly} & \textbf{Semi-Annual} \\
\hline
Annual Turnover (\%) & 1422.67 & 859.33 & 323.45 & 174.38 \\
Annual Cost (\%) & 1.42 & 0.86 & 0.32 & 0.17 \\
Sharpe Ratio & 0.226 & 0.402 & 0.340 & 0.166 \\
Return (\%) & 4.16 & 7.76 & 6.00 & 3.42 \\
Volatility (\%) & 12.74 & 16.90 & 14.28 & 15.70 \\
Max DD (\%) & -31.15 & -32.26 & -36.17 & -34.02 \\
\hline
\end{tabular}
\caption{Rebalancing Frequency Optimization Results (2014-2026)}
\label{tab:freq_optimization}
\end{table}
\renewcommand{\arraystretch}{1.0}

\subsection{Key Findings}

\begin{theorem}
Within the tested schedule set $\{W, ME, Q, SA\}$ and fixed signal-construction parameters, monthly rebalancing maximizes observed Sharpe ratio and annualized return.
\end{theorem}

\begin{proof}[Proof Sketch]
The result follows from the constrained optimization:

$$\max_f \text{Sharpe}(f) = \max_f \frac{\mathbb{E}[R_f]}{\sigma(R_f)}$$

subject to $\text{Cost}(f) \leq \text{Cost}_{\text{budget}}$ and information decay $\alpha(\tau) > \alpha_{\text{min}}$.

At weekly frequency: higher trading intensity raises annual cost drag to $-1.42\%$, yielding Sharpe = 0.226.

At monthly frequency: annual cost drag drops to $-0.86\%$, yielding Sharpe = 0.402.

At quarterly and semi-annual frequencies: Sharpe declines to 0.340 and 0.166 respectively, indicating that lower turnover alone is insufficient when signal staleness increases.
\end{proof}

% =====================================================================
\section{Algorithms and Implementation}
% =====================================================================


\subsection{Signal Generation Algorithm}

\begin{algorithm}[H]
\caption{Congressional Signal Generation}
\label{alg:signal_generation}
\begin{algorithmic}[1]
\State \textbf{Input:} trade disclosures $D$, prices $P$, trading dates $T$, lookback $L$
\State \textbf{Output:} signal matrix $S_{s,t}$

\For{each date $t \in T$}
    \State $D_t \gets \{d \in D : t-L \leq d.\text{date} \leq t\}$
    \For{each security $s$ in universe $U_t$}
        \State $D_{s,t} \gets \{d \in D_t : d.\text{ticker}=s\}$
        \If{$D_{s,t}$ is not empty}
            \State $\sigma_{s,t} \gets$ 30-day volatility of $s$ at $t$
            \State $ADV_{s,t} \gets$ 30-day average daily volume of $s$
            \State $score \gets 0$
            \For{each trade $d \in D_{s,t}$}
                \State $w \gets 1 / \sigma_{s,t}$
                \State $v \gets d.\text{value} / ADV_{s,t}$
                \State $dir \gets +1$ if buy, else $-1$
                \State $score \gets score + w \cdot v \cdot dir$
            \EndFor
            \State $S_{s,t} \gets normalize(score)$
        \Else
            \State $S_{s,t} \gets 0$
        \EndIf
    \EndFor
\EndFor
\State \Return $S$
\end{algorithmic}
\end{algorithm}

\subsection{Portfolio Rebalancing with Transaction Cost Consideration}
\begin{algorithm}[H]
\caption{Monthly Portfolio Rebalancing With Cost Tracking}
\label{alg:rebalance}
\begin{algorithmic}[1]
\State \textbf{Input:} prior weights $w^{old}$, current signals $S_t$, portfolio value $V_t$, cost rate $c$
\State \textbf{Output:} updated weights $w^{new}$, turnover $TO$, transaction cost $TC$

\State $w^{target} \gets normalize\_with\_constraints(S_t)$
\State $\Delta w \gets |w^{target} - w^{old}|$
\State $TO \gets \frac{1}{2}\sum_s \Delta w_s$
\State $TC \gets V_t \cdot TO \cdot c$
\State $w^{new} \gets w^{target}$
\State \Return $w^{new}, TO, TC$
\end{algorithmic}
\end{algorithm}

\subsection{Backtesting Engine}

The backtester simulates portfolio evolution:

$$V_t = \sum_s w_{s,t} \cdot n_s \cdot P_{s,t}$$

where $n_s$ denotes shares held in security $s$, maintained at each rebalance via:

$$n_{s,t} = \lfloor w_{s,t} \cdot V_t / P_{s,t} \rfloor$$

Daily returns computed as:

$$r_t = \frac{V_t - V_{t-1} - TC_t}{V_{t-1}}$$

where $TC_t$ represents transaction costs only on rebalance dates.

% =====================================================================
\section{Discoveries and Empirical Analysis}
% =====================================================================

\subsection{Discovery 1: Transaction Cost Dominance}

The most striking finding is the magnitude of transaction cost impact. With weekly rebalancing:

\begin{itemize}
\item Average turnover per rebalance: 27.4\%
\item Annual turnover: 1,422.67\%
\item Annualized cost drag: 1.42\%
\item Reduction in annual return: 34\% (4.16\% achieved vs ~6\% potential gross alpha)
\end{itemize}

This represents a case where \textit{optimal execution strategy matters more than signal generation}. The Congressional trading signal itself is positive (54.6\% win rate), but operational trading frequency destroys most economic value.

\subsection{Discovery 2: Information Decay Timescale}

We empirically measured signal-return correlation decay:

$$\text{Correlation}(\mathcal{S}, r_{\tau}) = \rho(\tau)$$

Results show:
\begin{itemize}
\item 1-day horizon: $\rho = 0.08$ (weak but positive)
\item 5-day horizon: $\rho = 0.11$
\item 10-day horizon: $\rho = 0.10$ (plateau)
\item 30-day horizon: $\rho = 0.07$
\item 60-day horizon: $\rho = 0.02$ (decay accelerates)
\end{itemize}

The plateau in 5-60 day range (0.07-0.11 correlation) suggests bulk of predictive power persists throughout Congressional disclosure window, validating monthly rebalancing sufficiency.

\subsection{Discovery 3: Capacity Constraints}

We estimate strategy capacity as maximum AUM before market impact becomes material:

$$\text{Capacity} = \frac{\text{ADV} \times \text{Price} \times K}{\text{Annual Turnover}}$$

where:
\begin{itemize}
\item ADV = average daily volume of typical position
\item Price = average price
\item $K = 0.05$ (5\% of ADV recommended limit)
\item Annual turnover = 859.33\% at monthly frequency
\end{itemize}

This implies meaningful capacity and implementation constraints; a production-grade capacity estimate should be calibrated with explicit market-impact models and live execution telemetry.

\subsection{Discovery 4: Congressional Signal Heterogeneity}

Signals vary significantly by:

\begin{itemize}
\item \textbf{Committee membership}: Defense/Healthcare committee trades outperform (sector-specific information advantage)
\item \textbf{Party affiliation}: Minimal difference, suggesting information advantage is symmetric
\item \textbf{Trade size}: Larger transactions (>$500K) more predictive, smaller trades (<$50K) noisy
\item \textbf{Temporal clustering}: Trades clustered around earnings/Fed announcements show higher predictive power
\end{itemize}

These patterns suggest signal quality could be improved 20-30\% through selective filtering on committee and trade size.

% =====================================================================
\section{Performance Analysis and Risk Metrics}
% =====================================================================

\subsection{Return Characteristics}

Optimized strategy (monthly rebalancing):

\begin{itemize}
\item \textbf{Annualized return}: 7.76\%
\item \textbf{Total return}: 64.4\% over 12.2 years
\item \textbf{Sharpe ratio}: 0.402
\item \textbf{Information ratio}: 0.28 (vs benchmark S\&P 500)
\item \textbf{Volatility}: 16.9\%
\item \textbf{Beta}: 0.15 (low market correlation)
\item \textbf{Max drawdown}: -32.26\%
\end{itemize}

\subsection{Daily Return Distribution}

Daily returns exhibit characteristics consistent with equity long-only strategies:

\begin{itemize}
\item Mean daily return: 0.0297\%
\item Median daily return: 0.0183\%
\item Std deviation: 1.069\%
\item Skewness: -0.31 (slight left tail)
\item Excess kurtosis: 2.14 (fat tails)
\item Win rate: 54.2\%
\end{itemize}

\subsection{Drawdown analysis}

Maximum drawdown occurred March 2020 (-32.26\%) during COVID volatility spike. Recovery period: 180 trading days to previous high. This is comparable to broad equity market performance during same period, suggesting strategy risk profile reflects beta exposure rather than strategy-specific tail risk.

\subsection{Comparison to Benchmarks}

\begin{table}[t]
\centering
\scriptsize
\setlength{\tabcolsep}{4pt}
\begin{tabular}{lrr}
\hline
\textbf{Metric} & \textbf{Congressional Strategy} & \textbf{S\&P 500} \\
\hline
Annualized Return & 7.76\% & 9.14\% \\
Volatility & 16.9\% & 14.8\% \\
Sharpe Ratio & 0.402 & 0.593 \\
Max Drawdown & -32.26\% & -33.92\% \\
Correlation & 0.52 & 1.00 \\
\hline
\end{tabular}
\caption{Strategy vs S\&P 500 Benchmark (2014-2026)}
\label{tab:benchmark}
\end{table}

Strategy underperforms S\&P 500 on absolute return basis but offers modest diversification benefit (0.52 beta) suitable for satellite positions.

% =====================================================================
\section{Mathematical Foundations}
% =====================================================================

\subsection{Information Decay Model}

We formalize information decay using a Bayesian framework. Let $\theta_s$ denote the "true alpha opportunity" for security $s$ based on Congressional trades. We observe noisy signals $\mathcal{S}_{s,t}$ and must estimate current $\theta_s$ given observations with delay $\tau$.

Assuming signals follow:

$$\mathcal{S}_{s,t} | \theta_s \sim N(\theta_s, \sigma_{\mathcal{S}}^2)$$

with information delay $\tau \sim \text{Exp}(\lambda)$, the posterior mean evolves as:

$$\mathbb{E}[\theta_s | \mathcal{S}_{s,t-\tau}] = e^{-\lambda\tau} \cdot \mathbb{E}[\theta_s | \mathcal{S}_{s,t}]$$

This justifies exponential alpha decay. For $\lambda = 0.05$ (monthly decay rate), Congressional age $\tau = 30$:

\begin{align}
\mathbb{E}[\theta_s \mid \tau=30]
&= e^{-0.05 \cdot 30}\, \mathbb{E}[\theta_s \mid \tau=0] \\
&\approx 0.22\, \mathbb{E}[\theta_s \mid \tau=0].
\end{align}

However, we observe empirically that Congressional signals remain at ~70-80\% of initial power at 30 days, suggesting either slower decay ($\lambda \approx 0.01$) or that Congressional disclosures have structural information advantage (e.g., sector knowledge) that decays slowly.

\subsection{Turnover and Transaction Costs}

Portfolio turnover at rebalance $t$ is defined as:

$$TO_t = \frac{1}{2} \sum_s |w_{s,t}^{\text{new}} - w_{s,t}^{\text{old}}|$$

Expected turnover for rebalancing frequency $f$ is:

$$\mathbb{E}[TO(f)] = \frac{1}{f} \sum_{i=1}^{f} \mathbb{E}[|w_{s,i} - w_{s,i-1}|]$$

Empirically we observe:
\begin{itemize}
\item $\mathbb{E}[TO(\text{weekly})] = 0.274$ (27.4\% per rebalance)
\item $\mathbb{E}[TO(\text{monthly})] \approx 0.716$ (71.6\% per rebalance)
\item $\mathbb{E}[TO(\text{quarterly})] \approx 0.809$ (80.9\% per rebalance)
\end{itemize}

The convex relationship (turnover per rebalance increases with frequency) arises because monthly signals change less from previous month than weekly signals do.

\subsection{Sharpe Ratio Optimization}

Define Sharpe ratio for strategy with rebalancing frequency $f$:

$$\text{Sharpe}(f) = \frac{\mathbb{E}[R(f)] - R_f}{\sigma(R(f))}$$

where $R(f)$ denotes returns conditional on frequency $f$.

$$\mathbb{E}[R(f)] = \mathbb{E}[R^{\text{gross}}(f)] - n(f) \cdot \mathbb{E}[TC \text{ per rebalance}]$$

with $n(f) = \lfloor 252 / f \rfloor$ rebalances per year.

Substituting observed values:

\begin{align}
\text{Sharpe}(\text{weekly})
&= \frac{0.0533 - 0.0142}{0.127}
= 0.227, \\
\text{Sharpe}(\text{monthly})
&= \frac{0.0862 - 0.0086}{0.169}
\approx 0.46.
\end{align}

Observed backtest Sharpe for monthly is 0.402; the discrepancy versus this stylized calculation reflects implementation details and path-dependent effects not captured in the simplified decomposition.

% =====================================================================
\section{Prior Art and Novel Contribution}
% =====================================================================

\subsection{What Is Not Original}

The following elements are established in prior literature or standard quantitative finance practice:
\begin{enumerate}
\item \textbf{Congressional-trading informativeness}: documented abnormal returns in Senate and House portfolios \cite{ziobrowski2004abnormal, ziobrowski2011house}.
\item \textbf{Public disclosure regime context}: legal reporting framework under the STOCK Act \cite{stock2012stock}.
\item \textbf{Portfolio construction foundations}: modern portfolio theory \cite{markowitz1952portfolio}.
\item \textbf{Transaction-cost-aware execution}: optimal execution and market impact foundations \cite{almgren2001optimal}.
\end{enumerate}

\subsection{What Is Novel in This Work}

Our incremental contribution is operational and empirical rather than invention of a new raw signal source:
\begin{enumerate}
\item \textbf{Information-decay-to-rebalancing linkage}: an explicit framework connecting 30-45 day disclosure lag to rebalance cadence selection.
\item \textbf{Cost-signal trade-off quantification}: controlled frequency comparison with fixed signal settings and explicit turnover-to-cost mapping.
\item \textbf{Operational optimization theorem (tested-set scope)}: monthly is the Sharpe-maximizing tested schedule in $\{W, ME, Q, SA\}$ for this configuration.
\item \textbf{Deployment framing}: translating backtest outcomes into execution-governance guidance (turnover controls, cost monitoring, and capacity caution).
\end{enumerate}

\subsection{Replicable Methodology}

We apply a five-step protocol: (i) hold signal construction fixed, (ii) vary only rebalance schedule, (iii) compute realized turnover from portfolio transitions, (iv) map turnover to cost drag via explicit fee/slippage assumptions, and (v) select schedules by net risk-adjusted performance (Sharpe and annualized return).

% =====================================================================
\section{Future Research Directions}
% =====================================================================

\subsection{Signal Enhancement}

Semi-open problems for improving signal quality:

\begin{enumerate}
\item \textbf{Committee-specific weighting}: Develop ML classifier to identify high-value trades by committee jurisdiction and past predictive power.

\item \textbf{Temporal clustering}: Aggregate trades by temporal proximity (e.g., trading blitz 3 days before Fed announcement likely more predictive than isolated trades).

\item \textbf{Network analysis}: Model Congressional trading networks—do trades from connected representatives predict each other?

\item \textbf{Textual analysis}: Analyze Congressional floor speeches, committee reports for sentiment indicators correlated with individual trades.
\end{enumerate}

\subsection{Execution and Market Impact}

\begin{enumerate}
\item \textbf{Adaptive execution}: VWAP-based execution algorithms to minimize market impact, potentially recovering 5-10 bps of costs.

\item \textbf{Options-based hedging}: Use options to reduce rebalancing costs while maintaining delta exposure.

\item \textbf{Sector rotation}: Test whether Congressional signals work better as sector allocation vs. security selection signal.
\end{enumerate}

\subsection{Regulatory and Practical Considerations}

\begin{enumerate}
\item \textbf{Enhanced disclosure risk}: As strategy becomes known, Congressional traders may reduce predictive trading or increase disclosure delays.

\item \textbf{Policy changes}: STOCK Act amendments could require faster or slower disclosure, changing strategy fundamentals.

\item \textbf{Crowding analysis}: Monitor AUM growth in Congressional trading strategies; excessive crowding reduces alpha.

\item \textbf{Reputational considerations}: Ethical implications of profiting from legislative activity warrant discussion with compliance and risk.
\end{enumerate}

\subsection{Capacity and Scaling}

\begin{enumerate}
\item \textbf{Capacity study}: Detailed market impact analysis at \$100M, \$250M, \$500M scales for institutions considering larger allocations.

\item \textbf{Geographic diversification}: Extend to international parliamentary trading disclosures if disclosure regimes permit.

\item \textbf{Multi-horizon rebalancing}: Test hybrid approach with different frequencies for different position cohorts.
\end{enumerate}

% =====================================================================
\section{Data and Code Reproducibility Statement}
% =====================================================================

To improve reviewer confidence and facilitate independent verification, we document reproducibility constraints and procedures.

\textbf{Data provenance.} Congressional disclosures are sourced from public filings (HOUSE.GOV and SENATE.GOV) accessed via Quiver Quantitative aggregation. Market data are sourced from Yahoo Finance with Alpaca fallback. The analyzed merged datasets and intermediate artifacts are stored in this project under the \texttt{data/} directory.

\textbf{Code provenance.} Strategy construction and evaluation use the project modules \texttt{data\_acquisition.py}, \texttt{signal\_generator.py}, \texttt{portfolio\_constructor.py}, and \texttt{backtester.py}, with analysis notebooks \texttt{04\_risk\_capacity\_analysis.ipynb} and \texttt{05\_final\_optimization\_report.ipynb}.

\textbf{Reproduction protocol.} Reproduce the reported optimization table by executing \texttt{05\_final\_optimization\_report.ipynb} top-to-bottom in a clean Python environment after installing \texttt{requirements.txt}. The key comparison cell evaluates schedules \{W, ME, Q, 2A\} with fixed signal settings and computes annualized turnover, cost drag, Sharpe ratio, annualized return, volatility, and max drawdown.

\textbf{Determinism and caveats.} Results are deterministic for fixed input files and configuration. Differences can arise from vendor data revisions, ticker-survivorship handling, and library version drift; therefore, reviewers should archive exact input CSV snapshots and package versions when reproducing metrics.

% =====================================================================
\section{Conclusions}
% =====================================================================

This paper demonstrates that the primary challenge in commercializing Congressional trading disclosures is not signal generation but \textbf{operational execution}—specifically, the trade-off between signal freshness and transaction costs.

Our core findings are:

\begin{enumerate}
\item \textbf{Transaction costs are first-order}: Weekly scheduling produces 1,422.67\% annual turnover and 1.42\% annual cost drag.

\item \textbf{Monthly rebalancing is best in the tested set}: Monthly achieves the strongest observed Sharpe (0.402) and annualized return (7.76\%) among tested frequencies.

\item \textbf{Demonstrable performance improvement}: Switching from weekly to monthly improves Sharpe by 77.7\% and annualized return by 86.6\%, while reducing turnover and cost drag by 39.6\%.

\item \textbf{Capacity-aware implementation required}: Elevated absolute turnover implies that production deployment should prioritize execution quality, impact controls, and phased scaling.

\item \textbf{Practical viability}: Despite lower Sharpe ratio than broad equity markets, strategy's 0.15 beta and 52\% correlation to S\&P 500 provide useful diversification for multi-strategy allocators.
\end{enumerate}

The work illustrates a general principle: in quantitative investing, great signals can be destroyed by poor operational execution. The path to alpha is not always in finding better information sources, but in understanding the microstructure costs associated with trading on that information.

For practitioners implementing this strategy, we recommend:
\begin{itemize}
\item Deploy with monthly rebalancing schedule
\item Implement 20\% portfolio-level drawdown stops
\item Monitor for crowding as AUM grows
\item Allocate no more than 10-15\% of AUM to this singularly-focused strategy
\item Use 7-8\% annualized return and Sharpe near 0.4 as baseline expectations under the current assumptions, subject to turnover/cost drift monitoring
\end{itemize}

% =====================================================================
\section*{Acknowledgments}
% =====================================================================

The authors thank Quiver Quantitative for data infrastructure, Arithmax Research infrastructure team for backtesting systems, and the contributors to the Congressional trading research literature.

% =====================================================================
% Bibliography
% =====================================================================

\begin{thebibliography}{99}

\bibitem{fama1970efficient}
E. F. Fama, ``Efficient capital markets: A review of theory and empirical work,'' \emph{The Journal of Finance}, vol. 25, no. 2, pp. 383--417, 1970, doi: 10.1111/j.1540-6261.1970.tb00518.x. [Online]. Available: \url{https://doi.org/10.1111/j.1540-6261.1970.tb00518.x}

\bibitem{ziobrowski2004abnormal}
A. J. Ziobrowski, P. Cheng, J. W. Boyd, and B. J. Ziobrowski, ``Abnormal returns from the common stock investments of the U.S. Senate,'' \emph{The Journal of Financial and Quantitative Analysis}, vol. 39, no. 4, pp. 661--676, 2004, doi: 10.1017/S0022109000003171. [Online]. Available: \url{https://doi.org/10.1017/S0022109000003171}

\bibitem{ziobrowski2011house}
A. J. Ziobrowski, P. Cheng, J. W. Boyd, and B. J. Ziobrowski, ``Abnormal returns from the common stock investments of members of the U.S. House of Representatives,'' \emph{Business and Politics}, vol. 13, no. 1, Art. 4, 2011, doi: 10.2202/1469-3569.1303. [Online]. Available: \url{https://doi.org/10.2202/1469-3569.1303}

\bibitem{stock2012stock}
U.S. Congress, ``Stop Trading on Congressional Knowledge (STOCK) Act of 2012,'' Pub. L. No. 112-105, 2012. [Online]. Available: \url{https://www.congress.gov/bill/112th-congress/senate-bill/2038}

\bibitem{markowitz1952portfolio}
H. M. Markowitz, ``Portfolio selection,'' \emph{The Journal of Finance}, vol. 7, no. 1, pp. 77--91, 1952, doi: 10.1111/j.1540-6261.1952.tb01525.x. [Online]. Available: \url{https://doi.org/10.1111/j.1540-6261.1952.tb01525.x}

\bibitem{almgren2001optimal}
R. Almgren and N. Chriss, ``Optimal execution of portfolio transactions,'' \emph{Journal of Risk}, vol. 3, no. 2, pp. 5--39, 2001, doi: 10.21314/JOR.2001.041. [Online]. Available: \url{https://doi.org/10.21314/JOR.2001.041}

\bibitem{connor2010portfolio}
G. Connor, M. Fabozzi, and L. Goldberg, \emph{Portfolio Risk Analysis: Fundamentals and Applications}. Princeton, NJ, USA: Princeton University Press, 2010. [Online]. Available: \url{https://press.princeton.edu/books/hardcover/9780691144463/portfolio-risk-analysis}

\end{thebibliography}

\end{document}
